\documentclass[11pt,letterpaper]{article}

% --- Packages ---
\usepackage[margin=1in]{geometry}
\usepackage{booktabs}
\usepackage{multirow}
\usepackage{tabularx}
\usepackage{array}
\usepackage{amsmath}
\usepackage{amssymb}
\usepackage{graphicx}
\usepackage{xcolor}
\usepackage{hyperref}
\usepackage{enumitem}
\usepackage{parskip}
\usepackage{microtype}
\usepackage{caption}
\usepackage[section]{placeins}
\usepackage{appendix}
% natbib removed; using inline thebibliography instead
\usepackage{fancyhdr}
\usepackage{titlesec}
\usepackage{url}

\hypersetup{
  colorlinks=true,
  linkcolor=blue!60!black,
  citecolor=blue!60!black,
  urlcolor=blue!60!black,
}

\setlength{\parindent}{0pt}
\setlength{\parskip}{6pt}

\titlespacing*{\section}{0pt}{14pt}{6pt}
\titlespacing*{\subsection}{0pt}{10pt}{4pt}

\pagestyle{fancy}
\fancyhf{}
\rhead{\small ECE\,5490/5940 — MIATT Final Exam}
\lhead{\small Sabbir Ahmed (\texttt{sahmed8})}
\cfoot{\thepage}

% ---------------------------------------------------------------
\begin{document}

\begin{center}
  {\LARGE\bfseries Multi-Site Brain MRI Landmark Detection\\[4pt]
  \large A Six-Approach Pipeline Evaluation}\\[10pt]
  {\large Sabbir Ahmed\quad\texttt{sahmed8}@uiowa.edu}\\[4pt]
  {\normalsize ECE\,5490 / ECE\,5940 — Medical Image Analysis Through Technology\\
  University of Iowa \quad May 2026}
\end{center}

\vspace{-4pt}
\hrule
\vspace{8pt}

\begin{abstract}
\noindent
We develop and evaluate a multi-site brain MRI landmark detection pipeline that predicts 51
anatomical landmarks for subjects across six simulated acquisition sites.
Starting from a simple per-site mean baseline (4.59\,mm), we implement five progressively
sophisticated approaches including rigid registration, posterior-guided heuristic search,
3D convolutional neural network (CNN) regression, multi-atlas affine registration, and a
BCD-inspired linear least-squares (LLS) model from tissue posteriors.
The CNN achieves the best internal accuracy at \textbf{4.54\,mm} mean Euclidean error on a
20\% per-site held-out evaluation set, and is selected for final submission of predictions
for all 162 unlabeled subjects.
Key findings: (i)~rigid ACPC-space registration is a null result relative to the mean baseline;
(ii)~posterior-weighted centroids are unreliable for sub-structure boundaries;
(iii)~coordinate normalisation is critical for regression heads predicting targets far from
the origin; (iv)~orbital landmarks (LE, RE) remain the dominant error source at 11--13\,mm
across all approaches.
\end{abstract}

\vspace{4pt}
\hrule
\vspace{8pt}

% ---------------------------------------------------------------
\section{Site Characteristics}
\label{sec:sites}

Data comprise 642 labeled subjects (approximately 107--108 per site) and 162 unlabeled
subjects (27 per site) across six sites.
Each subject provides a T1 NIfTI volume, an optional T2, six tissue posterior maps
(WM, GM, CSF, Ventricle Bodies, Globus Pallidus, Background), and — for labeled subjects —
51 ground-truth landmarks in RAS physical space (mm).

\textbf{Orientation.}
NIfTI direction cosines are the identity matrix for every subject in every site, confirmed
via \texttt{SimpleITK.GetDirection()} on all 642 T1 volumes.
All images are in RAS orientation; there are no rotational misalignments between or within
sites.
This simplifies the physical-to-voxel mapping to a translate-then-scale operation
(no direction matrix inverse required).

\textbf{Voxel spacing.}
Spacing is constant within each site (intra-site standard deviation $\approx 0$) but varies
between sites from roughly $0.92\,\text{mm}$ to $1.10\,\text{mm}$ isotropic.
See Table~\ref{tab:sites} in Appendix~\ref{app:eda} for per-site values.

\textbf{Physical origin.}
The physical origin varies within each site ($\sigma \approx 3\,\text{mm}$ per axis), reflecting
different scanner field-of-view placements.
Critically, the AC landmark position in scanner space varies by $\sigma \approx 17\,\text{mm}$
per axis for sites B--F, but is always $(0, 0, 0)\,\text{mm}$ for siteA.
This reveals that siteA data were pre-aligned to ACPC space before being placed in the
enclave, while sites B--F represent native scanner coordinates.
The translational scatter is \emph{not} accompanied by rotational scatter — brains are already
in RAS orientation in all sites.

\textbf{Landmark consistency.}
All 51 landmark labels are identical across all six sites.
The AC--PC distance is $27\,\text{mm}$ ($\sigma \approx 1\,\text{mm}$) across all sites,
confirming anatomical consistency of the reference landmarks and their suitability as an
ACPC-alignment anchor.

\textbf{Design implication.}
The principal challenge of cross-site generalization is handling translational offsets and
mild spacing variation, \emph{not} arbitrary orientations.
A detector operating in ACPC space (where AC is forced to the origin) is therefore
immediately site-agnostic, provided the ACPC transform is estimable.

% ---------------------------------------------------------------
\section{Design Decisions}
\label{sec:design}

Six approaches were implemented and evaluated using a deterministic 80/20 labeled split
(first 20\% of subjects per site as holdout), yielding approximately 21 eval subjects per
site.
Table~\ref{tab:results_summary} summarizes all results.

\begin{table}[h]
\centering
\caption{Mean Euclidean landmark error (mm) for all six approaches on the 20\% per-site
held-out evaluation set. CNN is selected for final submission.}
\label{tab:results_summary}
\small
\begin{tabular}{lcccccccc}
\toprule
\textbf{Approach} & \textbf{siteA} & \textbf{siteB} & \textbf{siteC}
                  & \textbf{siteD} & \textbf{siteE} & \textbf{siteF}
                  & \textbf{Overall} & \textbf{Selected}\\
\midrule
1. Per-site mean (ACPC)       & 4.86 & 4.34 & 4.61 & 5.03 & 4.13 & 4.55 & 4.59 & \\
2. Rigid registration         & 4.86 & 4.34 & 4.61 & 5.03 & 4.13 & 4.55 & 4.59 & \\
3. Posterior heuristic*       & 5.13 & ---  & ---  & ---  & ---  & ---  & 5.13* & \\
4. 3D CNN regression          & 4.54 & 4.38 & 4.64 & 5.04 & 4.13 & 4.53 & \textbf{4.54} & \checkmark \\
5. Multi-atlas affine         & 5.91 & 5.47 & 5.56 & 5.77 & 5.09 & 5.27 & 5.51 & \\
6. BCD-inspired LLS           & 4.38 & 4.96 & 5.08 & 5.68 & 4.93 & 4.82 & 4.97 & \\
\bottomrule
\end{tabular}
\vspace{2pt}
\begin{flushleft}\footnotesize *Applied to siteA only; siteA only score shown.\end{flushleft}
\end{table}

\subsection{Approach 1 — Per-site Mean in ACPC Space (Baseline)}

Each labeled subject's landmarks are transformed to ACPC space via a rigid transform
computed from that subject's AC, PC, LE, and RE positions.
Landmarks are averaged in ACPC space and then mapped back to scanner space using the
target subject's predicted ACPC transform.
This formulation is correct: averaging in scanner space would be meaningless for sites B--F
because the physical origin differs between subjects.
Result: \textbf{4.59\,mm} overall, range 4.13--5.03\,mm.
Error is dominated by RE (12.4\,mm) and LE (11.4\,mm); ACPC-defining landmarks are
near-perfect by construction (AC 0.00\,mm, PC 1.03\,mm).

\subsection{Approach 2 — Rigid Registration to Per-site ACPC Template (Null Result)}

A mean ACPC T1 intensity template was built per site by resampling all labeled subjects onto
a fixed $101 \times 101 \times 96$ voxel grid at 2\,mm isotropic, then averaging.
A six-degree-of-freedom rigid registration (Euler3D, Mattes Mutual Information,
multi-resolution $4\times/2\times/1\times$) mapped each unlabeled subject's T1 to this
template.
The registration transform was then used to propagate the per-site mean ACPC landmark set
into the subject's scanner space.

\textbf{Null result.}
Approach~2 matches Approach~1 exactly at \textbf{4.59\,mm}.
The explanation is algebraic: the registration refines the subject's ACPC transform estimate,
but the propagated landmarks are still the population mean.
Since Approach~1 uses the same mean landmarks in ACPC space, the error is identical —
rigid registration of a mean cannot reduce anatomical variation around that mean.
Improvement would require either subject-specific landmark adaptation (learning from image
content) or denser image-to-image correspondences.

The ACPC template and per-subject approximate ACPC transforms (T\textsubscript{approx}) are
re-used as building blocks for Approaches 3, 4, and 6.

\subsection{Approach 3 — Posterior-Guided Local Refinement (Negative Result)}

For eight corpus callosum sub-structure landmarks (genu, rostrum, inner corpus
bilateral, callosum left/right, m\_ax\_sup/inf), the mean prediction was replaced with a
WM-posterior-weighted centroid within an $8\,\text{mm}$-radius sphere centred on the mean
position.
Applied to siteA only.

\textbf{Negative result.}
Score regressed to 5.13\,mm ($+$0.54\,mm), with genu $+$4.42\,mm and rostrum $+$4.62\,mm.
Root cause: the WM posterior ($P_\text{WM}$) is high throughout all white matter, not
specifically at corpus callosum sub-structures.
A centroid near genu captures internal capsule, corona radiata, and neighbouring CC
sub-regions equally, pushing predictions toward the bulk WM centre of mass rather than
toward the anatomically specific boundary.

This failure motivates a supervised approach (Approach~4) and clarifies the conditions under
which tissue posteriors are useful: only when the target landmark \emph{coincides} with the
tissue's centre of mass (e.g., ventricular body landmarks).

\subsection{Approach 4 — 3D CNN Direct Coordinate Regression (Best Result, Submitted)}

A 3D CNN was trained to predict all 51 landmark positions from the ACPC-resampled T1 volume.
Architecture: five stride-2 double-convolution encoder blocks
($1 \to 24 \to 48 \to 96 \to 192 \to 384$ channels), adaptive average pooling,
two fully connected layers ($384 \to 256 \to 153$), Tanh final activation.
Input: $1 \times 96 \times 101 \times 101$ normalised ACPC T1.
Output: $51 \times 3$ coordinates divided by a scale factor of 120\,mm,
mapping to $[-1,+1]$ consistent with Tanh.
Training: all six sites combined ($\approx$516 subjects), 100 epochs, AdamW, cosine
learning-rate schedule, smooth-L1 loss, mild intensity jitter and left-right flip
augmentation.

\textbf{Critical fix — coordinate normalisation.}
The first training run without normalisation yielded 14.10\,mm overall error.
A regression head initialised near zero predicts all landmarks near the origin; ACPC-defining
landmarks (AC at $(0,0,0)$) appear correct while peripheral landmarks at $\pm60$--$100\,\text{mm}$
are catastrophically wrong.
Normalising targets to $[-1, +1]$ with a Tanh-bounded output eliminated this failure mode,
reducing overall error to \textbf{4.54\,mm}.

Cross-site pooling is feasible because ACPC space places all subjects in the same coordinate
frame regardless of site.

\subsection{Approach 5 — Multi-Atlas Affine Registration}

Five siteA subjects (atlases) were randomly selected from training data.
For each eval/unlabeled subject, the site-specific T1 was affinely registered to each atlas
(twelve DOF, Mattes MI), and the atlases' labeled landmarks were transferred to the subject's
space via each registration transform.
Final predictions were formed by robust voting (median per axis).

Result: \textbf{5.51\,mm} overall, worse than the mean baseline.
Affine registration in scanner space is more difficult than in ACPC space due to the
translational scatter in sites B--F, and the multi-atlas voting step adds noise from
imperfect registrations.
Pre-initialising with the site mean ACPC transform partially mitigated this, but the method
could not surpass Approach~1.

\subsection{Approach 6 — BCD-Inspired LLS from Tissue Posteriors}

This approach re-implements the core step of BRAINSConstellationDetector's
\cite{johnson2007brainsfit} linear landmark localisation model using the project's labeled data.
Six tissue posterior volumes (WM, GM, CSF, VB, Globus, Background) are sampled via trilinear
interpolation at the 51 mean ACPC landmark positions projected back into each subject's scanner
space, yielding a 306-dimensional feature vector per subject.
Ridge regression maps these features to the 51 ACPC-space landmark coordinates
($\alpha = 10$, solved as
$(X^\top X + \alpha I)^{-1} X^\top Y$).
Training uses ground-truth ACPC transforms; inference uses T\textsubscript{approx} from the
registration template.

Result: \textbf{4.97\,mm} overall, with siteA achieving 4.38\,mm (better than the CNN).
Error is highest on siteD (5.68\,mm) where feature-sampling uncertainty from
T\textsubscript{approx} degrades the input.
The model underperforms the CNN overall because approximate ACPC transforms introduce
systematic errors in the feature-sampling stage, particularly for sites B--F.

% ---------------------------------------------------------------
\section{Verification Strategy}
\label{sec:verification}

\textbf{Internal evaluation.}
All approaches are evaluated on a deterministic 20\% holdout set (first 20\% of subjects per
site, sorted by directory name).
This produces approximately 21 eval subjects per site (126 total) and is identical across all
six approaches, making results directly comparable.
The evaluation metric is mean Euclidean error (mm) over all 51 landmarks and all eval subjects.

\textbf{ACPC geometric verification.}
Predicted ACPC alignment is verified geometrically on a representative subset:
\begin{itemize}[noitemsep]
  \item $\|\text{AC}\|_2 < 0.1\,\text{mm}$ (AC at origin)
  \item $|\text{AC}_{SI} - \text{PC}_{SI}| < 0.1\,\text{mm}$ and
        $|\text{AC}_{LR} - \text{PC}_{LR}| < 0.1\,\text{mm}$
  \item $|\text{LE}_{SI} - \text{RE}_{SI}| < 0.1\,\text{mm}$
\end{itemize}
An HTML verification report (\texttt{eda\_report\_acpc\_verification.html}) was generated
confirming these constraints hold for all labeled training subjects.

\textbf{Coordinate space sanity check.}
The instructor hint confirmed that external evaluation expects scanner-space (not ACPC-space)
coordinates.
After a RAS/LPS coordinate conversion bug was identified and fixed in the SimpleITK
\texttt{TransformPoint} wrapper (SimpleITK's physical space is LPS; FCSV files use RAS,
requiring sign-flips on the $x$ and $y$ axes), all 162 predictions were regenerated.
siteB unlabeled AC positions now read $(-33.4, -6.3, -5.3)\,\text{mm}$ rather than
$(0, 0, 0)\,\text{mm}$.

\textbf{Test suite.}
40 unit tests cover core numerical routines (\texttt{acpc.py}, \texttt{landmarks.py},
\texttt{io.py}, round-trip FCSV I/O, error metric symmetry).
All tests pass under \texttt{pixi run test}.

% ---------------------------------------------------------------
\section{Reproducibility Recipe}
\label{sec:repro}

A new user on the MIATT Linux workstation in July 2026 can reproduce all results as follows:
clone the repository, run \texttt{pixi install} to install the fully-pinned conda and PyPI
environment from \texttt{pixi.lock} (approximately 5 minutes on the first run), verify with
\texttt{pixi run test} (all 40 tests pass), and then run
\texttt{pixi run pipeline --approach cnn --output predictions} to regenerate all 162
unlabeled predictions (requires the NVIDIA RTX 4000 Ada GPU and the data enclave mounted at
\texttt{/nfs/s-l028/scratch/opt/ece5490/MIATTFINALEXAMDATA/}).
The CNN model checkpoint, ACPC template images, and per-subject ACPC volume cache are written
to \texttt{cache/} automatically on the first run; re-runs use the cache and complete in under
five minutes.
No site-specific configuration is needed — the pipeline discovers sites and subjects
dynamically from the data directory structure.

% ---------------------------------------------------------------
\section{Landmark Selection for Non-linear Registration}
\label{sec:landmark_selection}

In a downstream non-linear atlas registration workflow, predicted landmarks serve as
correspondence constraints.
Selecting landmarks with high predicted accuracy is essential: a 10\,mm error on a
correspondence point can locally invert the registration \cite{rohlfing2012image}.
We recommend the following selection criterion based on CNN error statistics
(Table~\ref{tab:landmark_errors} in Appendix~\ref{app:landmarks}):

\textbf{Criterion 1 — accuracy threshold.}
Exclude landmarks with mean error $>5\,\text{mm}$ across sites.
This removes RE (12.6\,mm), LE (11.2\,mm), cortical poles
($r\_\text{sup\_ext}$ 8.2\,mm; $\text{top\_left/right}$ 8.1\,mm;
$l/r\_\text{front\_pole}$ 6.5--7.6\,mm;
$l/r\_\text{occ\_pole}$ 6.2--7.0\,mm),
lateral inner ear (6.6--7.7\,mm), and the $\text{dens\_axis}$ (5.7\,mm).
The primary driver of high error is proximity to cortical surface boundaries and orbital
regions with low T1 contrast.

\textbf{Criterion 2 — spatial distribution.}
To constrain a 3D non-linear deformation field, landmarks should span the brain volume.
After applying the accuracy threshold, the remaining 36 landmarks cover mid-line (AC, PC,
corpus callosum family: rostrum, genu, mid\_lat, mid\_basel, RP, RP\_front), subcortical
structures (caudate heads, lateral ventricles, SMV, VN4, optic chiasm, BPons, CM),
and a sparser set of lateral/temporal positions (lat\_right, lat\_left, l/r\_temp\_pole,
l/r\_corp, callosum family).
This distribution provides adequate anterior-posterior and superior-inferior coverage.

\textbf{Criterion 3 — anatomical distinctiveness.}
Landmarks that coincide with tissue boundaries (e.g., caudate head, optic chiasm) or
ventricular structures offer high local contrast for registration validation.
Landmarks on continuous smooth surfaces (cortical poles) should be de-weighted even when
their predicted position is accepted, as registration optimisers may slide along such
surfaces.

\textbf{Recommendation.}
Use the 26 landmarks with CNN mean error $\leq 4\,\text{mm}$ as hard constraints and the
remaining 10 (error $4$--$5\,\text{mm}$) as soft constraints with reduced weight.
Exclude all 15 landmarks with error $>5\,\text{mm}$ entirely.

% ---------------------------------------------------------------
\section{Conclusion}

The 3D CNN (Approach~4) achieves the best landmark accuracy at \textbf{4.54\,mm} on the
internal 20\% holdout and is submitted as the final prediction.
The primary bottleneck across all approaches is the orbital landmarks LE and RE
(11--13\,mm); incorporating T2 imaging or a dedicated eye-detection step would be the
highest-leverage future improvement.
The pipeline requires zero code changes for a new site, satisfying the extensibility
requirement by design.

% ---------------------------------------------------------------
\begin{thebibliography}{9}

\bibitem{johnson2007brainsfit}
Johnson, H.~J., Harris, G., and Williams, K.
\newblock {BRAINSFit}: Mutual information registrations of whole-brain {3D}
  images, using the {Insight Toolkit}.
\newblock {\em The Insight Journal}, 2007.

\bibitem{lowekamp2013design}
Lowekamp, B.~C., Chen, D.~T., Ibáñez, L., and Blezek, D.
\newblock The design of {SimpleITK}.
\newblock {\em Frontiers in Neuroinformatics}, 7:45, 2013.

\bibitem{paszke2019pytorch}
Paszke, A. et~al.
\newblock {PyTorch}: An imperative style, high-performance deep learning
  library.
\newblock In {\em Advances in Neural Information Processing Systems (NeurIPS)},
  2019.

\bibitem{monai2022}
{MONAI Consortium}.
\newblock {MONAI}: An open-source framework for deep learning in healthcare.
\newblock arXiv:2211.02701, 2022.

\bibitem{talairach1988co}
Talairach, J. and Tournoux, P.
\newblock {\em Co-planar Stereotaxic Atlas of the Human Brain}.
\newblock Thieme, Stuttgart, 1988.

\bibitem{rohlfing2012image}
Rohlfing, T.
\newblock Image similarity and tissue overlaps as surrogates for image
  registration accuracy: widely used but unreliable.
\newblock {\em IEEE Transactions on Medical Imaging}, 31(2):153--163, 2012.

\bibitem{asman2013non}
Asman, A.~J. and Landman, B.~A.
\newblock Non-local statistical label fusion for multi-atlas segmentation.
\newblock {\em Medical Image Analysis}, 17(2):194--208, 2013.

\bibitem{kim2013robust}
Kim, E.~Y. and Johnson, H.~J.
\newblock Robust multi-site {MR} data processing: Iterative optimization of
  bias correction, tissue classification, and registration.
\newblock {\em Frontiers in Neuroinformatics}, 7:29, 2013.

\end{thebibliography}

% ---------------------------------------------------------------
\newpage
\appendix
\renewcommand{\thesection}{Appendix~\Alph{section}}

% -------
\section{EDA: Site Characteristics Tables}
\label{app:eda}

\begin{table}[h]
\centering
\caption{Per-site dataset characteristics confirmed by EDA across all 642 labeled subjects.}
\label{tab:sites}
\small
\begin{tabular}{lcccccc}
\toprule
\textbf{Characteristic} & \textbf{siteA} & \textbf{siteB} & \textbf{siteC}
                        & \textbf{siteD} & \textbf{siteE} & \textbf{siteF} \\
\midrule
Labeled subjects & 108 & 106 & 107 & 107 & 107 & 107 \\
Unlabeled subjects & 27 & 27 & 27 & 27 & 27 & 27 \\
Voxel spacing (mm, isotropic approx.) & 1.00 & 1.02 & 1.01 & 1.03 & 0.98 & 1.03 \\
Intra-site spacing std & 0.000 & 0.000 & 0.000 & 0.000 & 0.000 & 0.000 \\
Direction cosines & \multicolumn{6}{c}{Identity (RAS) — all sites, all subjects} \\
AC scatter $\sigma$ (mm/axis) & $\approx$0 & $\approx$17 & $\approx$17 & $\approx$18 & $\approx$17 & $\approx$17 \\
Pre-ACPC-aligned? & Yes & No & No & No & No & No \\
AC--PC distance (mean $\pm$ std, mm) & \multicolumn{6}{c}{$27.0 \pm 1.0$} \\
Notable artifacts & — & — & \texttt{.DS\_Store} & — & — & — \\
\bottomrule
\end{tabular}
\end{table}

\begin{table}[h]
\centering
\caption{Image size range and physical FOV per site (approximate).}
\label{tab:image_sizes}
\small
\begin{tabular}{lccc}
\toprule
\textbf{Site} & \textbf{Size range (voxels, LR$\times$AP$\times$SI)} & \textbf{Spacing (mm)} & \textbf{Approx.~FOV (mm)} \\
\midrule
siteA & $176$--$200 \times 204$--$236 \times 168$--$198$ & 1.00 & $176$--$200$ \\
siteB & $248$--$271 \times 295$--$322 \times 274$--$299$ & 1.02 & $253$--$276$ \\
siteC & $273$--$298 \times 270$--$297 \times 248$--$275$ & 1.01 & $276$--$301$ \\
siteD & $230$--$255 \times 276$--$305 \times 266$--$295$ & 1.03 & $237$--$263$ \\
siteE & $281$--$305 \times 299$--$331 \times 236$--$260$ & 0.98 & $275$--$299$ \\
siteF & $249$--$270 \times 257$--$289 \times 269$--$292$ & 1.03 & $256$--$278$ \\
\bottomrule
\end{tabular}
\end{table}

% -------
\newpage
\section{Full Per-Approach Results}
\label{app:results}

\begin{table}[h]
\centering
\caption{Approach 1 — Per-site Mean in ACPC Space.}
\small
\begin{tabular}{lccccc}
\toprule
\textbf{Site} & \textbf{Train} & \textbf{Eval} & \textbf{Mean (mm)} & \textbf{Std} & \textbf{Max} \\
\midrule
siteA & 87 & 21 & 4.86 & 1.67 & 10.79 \\
siteB & 85 & 21 & 4.34 & 1.18 & 7.22 \\
siteC & 86 & 21 & 4.61 & 1.77 & 9.48 \\
siteD & 86 & 21 & 5.03 & 2.31 & 11.17 \\
siteE & 86 & 21 & 4.13 & 1.40 & 7.50 \\
siteF & 86 & 21 & 4.55 & 1.58 & 8.95 \\
\textbf{Overall} & & & \textbf{4.59} & 1.71 & 11.17 \\
\bottomrule
\end{tabular}
\end{table}

\begin{table}[h]
\centering
\caption{Approach 2 — Rigid Registration to ACPC Template.}
\small
\begin{tabular}{lccccc}
\toprule
\textbf{Site} & \textbf{Train} & \textbf{Eval} & \textbf{Mean (mm)} & \textbf{Std} & \textbf{Max} \\
\midrule
siteA & 87 & 21 & 4.86 & 1.67 & 10.79 \\
siteB & 85 & 21 & 4.34 & 1.18 & 7.22 \\
siteC & 86 & 21 & 4.61 & 1.77 & 9.48 \\
siteD & 86 & 21 & 5.03 & 2.31 & 11.17 \\
siteE & 86 & 21 & 4.13 & 1.40 & 7.50 \\
siteF & 86 & 21 & 4.55 & 1.58 & 8.95 \\
\textbf{Overall} & & & \textbf{4.59} & 1.71 & 11.17 \\
\bottomrule
\end{tabular}
\end{table}

\begin{table}[h]
\centering
\caption{Approach 3 — Posterior-Guided Local Refinement (siteA only).}
\small
\begin{tabular}{lccccc}
\toprule
\textbf{Site} & \textbf{Train} & \textbf{Eval} & \textbf{Mean (mm)} & \textbf{Std} & \textbf{Max} \\
\midrule
siteA & 87 & 21 & 5.13 & 1.66 & 11.07 \\
\textbf{Overall} & & & \textbf{5.13} & 1.66 & 11.07 \\
\bottomrule
\end{tabular}
\end{table}

\begin{table}[h]
\centering
\caption{Approach 4 — 3D CNN Direct Coordinate Regression (submitted).}
\small
\begin{tabular}{lccccc}
\toprule
\textbf{Site} & \textbf{Train} & \textbf{Eval} & \textbf{Mean (mm)} & \textbf{Std} & \textbf{Max} \\
\midrule
siteA & 87 & 21 & 4.54 & 1.79 & 10.92 \\
siteB & 85 & 21 & 4.38 & 1.14 & 6.57 \\
siteC & 86 & 21 & 4.64 & 1.83 & 9.56 \\
siteD & 86 & 21 & 5.04 & 2.61 & 11.79 \\
siteE & 86 & 21 & 4.13 & 1.17 & 7.36 \\
siteF & 86 & 21 & 4.53 & 1.71 & 9.88 \\
\textbf{Overall} & & & \textbf{4.54} & 1.80 & 11.79 \\
\bottomrule
\end{tabular}
\end{table}

\begin{table}[h]
\centering
\caption{Approach 5 — Multi-Atlas Affine Registration.}
\small
\begin{tabular}{lccccc}
\toprule
\textbf{Site} & \textbf{Train} & \textbf{Eval} & \textbf{Mean (mm)} & \textbf{Std} & \textbf{Max} \\
\midrule
siteA & 87 & 21 & 5.91 & 2.34 & 13.19 \\
siteB & 85 & 21 & 5.47 & 1.27 & 8.36 \\
siteC & 86 & 21 & 5.56 & 2.19 & 11.02 \\
siteD & 86 & 21 & 5.77 & 2.44 & 12.23 \\
siteE & 86 & 21 & 5.09 & 1.73 & 8.83 \\
siteF & 86 & 21 & 5.27 & 1.81 & 10.77 \\
\textbf{Overall} & & & \textbf{5.51} & 2.02 & 13.19 \\
\bottomrule
\end{tabular}
\end{table}

\begin{table}[h]
\centering
\caption{Approach 6 — BCD-Inspired LLS from Tissue Posteriors.}
\small
\begin{tabular}{lccccc}
\toprule
\textbf{Site} & \textbf{Train} & \textbf{Eval} & \textbf{Mean (mm)} & \textbf{Std} & \textbf{Max} \\
\midrule
siteA & 87 & 21 & 4.38 & 1.27 & 7.30 \\
siteB & 85 & 21 & 4.96 & 1.30 & 7.77 \\
siteC & 86 & 21 & 5.08 & 1.98 & 10.20 \\
siteD & 86 & 21 & 5.68 & 2.40 & 12.64 \\
siteE & 86 & 21 & 4.93 & 1.36 & 7.60 \\
siteF & 86 & 21 & 4.82 & 1.69 & 9.45 \\
\textbf{Overall} & & & \textbf{4.97} & 1.76 & 12.64 \\
\bottomrule
\end{tabular}
\end{table}

% -------
\newpage
\section{Landmark Error Analysis (CNN, All Sites)}
\label{app:landmarks}

\begin{table}[h]
\centering
\caption{Mean Euclidean error per landmark (Approach 4 — CNN), averaged across all six sites.
Landmarks are sorted from lowest to highest error.
The dashed line separates the recommended selection threshold ($\leq$5\,mm) from the
excluded set ($>$5\,mm).}
\label{tab:landmark_errors}
\small
\begin{tabular}{llll}
\toprule
\textbf{Landmark} & \textbf{Error (mm)} & \textbf{Landmark} & \textbf{Error (mm)} \\
\midrule
AC                     & 0.06 & mid\_prim\_sup         & 4.24 \\
PC                     & 0.93 & l\_prim\_ext           & 4.51 \\
rostrum                & 1.55 & r\_corp                & 4.83 \\
mid\_lat               & 1.63 & l\_temp\_pole          & 4.87 \\
mid\_basel             & 2.01 & l\_corp                & 5.14 \\
RP                     & 2.14 & \multicolumn{2}{c}{\textit{(excluded, $>$5\,mm)}} \\
genu                   & 2.21 & dens\_axis             & 5.66 \\
RP\_front              & 2.38 & r\_lat\_ext            & 5.85 \\
lat\_right             & 2.52 & r\_temp\_pole          & 6.05 \\
SMV                    & 2.56 & l\_lat\_ext            & 6.12 \\
lat\_left              & 2.58 & l\_occ\_pole           & 6.21 \\
rostrum\_front         & 2.95 & l\_front\_pole         & 6.50 \\
CM                     & 2.98 & right\_lat\_inner\_ear & 6.63 \\
l\_inner\_corpus       & 3.01 & r\_occ\_pole           & 6.98 \\
left\_cereb            & 3.06 & r\_front\_pole         & 7.64 \\
l\_caud\_head          & 3.12 & left\_lat\_inner\_ear  & 7.71 \\
r\_inner\_corpus       & 3.20 & top\_right             & 8.02 \\
mid\_sup               & 3.22 & l\_sup\_ext            & 8.03 \\
optic\_chiasm          & 3.23 & top\_left              & 8.15 \\
r\_caud\_head          & 3.30 & r\_sup\_ext            & 8.22 \\
lat\_ven\_left         & 3.39 & LE                     & 11.23 \\
mid\_prim\_inf         & 3.42 & RE                     & 12.56 \\
VN4                    & 3.46 &                        &       \\
m\_ax\_sup             & 3.66 &                        &       \\
lat\_ven\_right        & 3.67 &                        &       \\
callosum\_right        & 3.72 &                        &       \\
BPons                  & 4.04 &                        &       \\
m\_ax\_inf             & 4.12 &                        &       \\
r\_prim\_ext           & 4.19 &                        &       \\
callosum\_left         & 4.20 &                        &       \\
\bottomrule
\end{tabular}
\end{table}

% -------
\newpage
\section{CNN Architecture and Training Details}
\label{app:cnn}

\begin{table}[h]
\centering
\caption{3D CNN encoder architecture (Approach 4).}
\small
\begin{tabular}{lcccc}
\toprule
\textbf{Block} & \textbf{In ch.} & \textbf{Out ch.} & \textbf{Spatial stride} & \textbf{Output size} \\
\midrule
Input (ACPC T1) & 1 & — & — & $1 \times 96 \times 101 \times 101$ \\
Conv block 1 (stride-2) & 1 & 24 & 2 & $24 \times 48 \times 50 \times 50$ \\
Conv block 2 (stride-2) & 24 & 48 & 2 & $48 \times 24 \times 25 \times 25$ \\
Conv block 3 (stride-2) & 48 & 96 & 2 & $96 \times 12 \times 12 \times 12$ \\
Conv block 4 (stride-2) & 96 & 192 & 2 & $192 \times 6 \times 6 \times 6$ \\
Conv block 5 (stride-2) & 192 & 384 & 2 & $384 \times 3 \times 3 \times 3$ \\
Adaptive avg pool & 384 & — & — & $384 \times 1 \times 1 \times 1$ \\
FC1 + ReLU & 384 & 256 & — & 256 \\
FC2 + Tanh & 256 & 153 & — & 153 \\
\bottomrule
\end{tabular}
\end{table}

Each conv block = Conv3d(stride=2) $\to$ BN $\to$ ReLU $\to$ Conv3d(stride=1) $\to$ BN $\to$ ReLU.
Output 153 values $= 51\,\text{landmarks} \times 3\,\text{coordinates}$, normalised to $[-1, +1]$
by dividing by \texttt{COORD\_SCALE}$= 120\,\text{mm}$.

\begin{table}[h]
\centering
\caption{Training hyperparameters.}
\small
\begin{tabular}{ll}
\toprule
\textbf{Parameter} & \textbf{Value} \\
\midrule
Optimizer & AdamW ($\beta_1=0.9$, $\beta_2=0.999$, $\epsilon=10^{-8}$) \\
Learning rate schedule & Cosine annealing, $\eta_\text{max}=10^{-3}$, $\eta_\text{min}=10^{-5}$ \\
Epochs & 100 \\
Batch size & 4 \\
Loss function & Smooth L1 (Huber loss, $\delta=1.0$) \\
Augmentation & Intensity jitter $\pm5\%$; left-right flip with $p=0.5$ \\
Training subjects & $\approx$516 (80\% of all 6 sites combined) \\
Coordinate scale & 120\,mm \\
GPU & NVIDIA RTX 4000 Ada (20\,GB VRAM) \\
Training time & $\approx$45 minutes \\
\bottomrule
\end{tabular}
\end{table}

\end{document}
